\section{Proofs for Reflection, Ballot Counts and Arcsine Laws}
\label{sec:annex-p15-reflection-and-ballot-counts}

\begin{annexitem}{annex:p15.rb-reflection}{thm:p15-rb-reflection}{The exact first-zero involution}
\begin{proof}
A path from zero to \(z\) has \(a\) up-steps and \(b\) down-steps
with \(a+b=n\), \(a-b=z\). Thus \(a=(n+z)/2\) and
\(b=(n-z)/2\). These must be nonnegative integers. Choosing
the \(a\) up-step positions gives \(\binom na\) paths.
If either requirement fails there is no path, in agreement
with the stated zero convention. Independent fair increments
give probability \(2^{-n}\) to every such path.
Translation gives \(N_{n,y-x}\) unrestricted paths from \(x\)
to \(y\).

Now take positive integers \(x,y\) and a path \(s\) from \(x\)
to \(y\) that visits zero. Let
\(K=\min\{0\leq j\leq n:s_j=0\}\).
Since the endpoints are positive, \(1\leq K<n\).
Define
\[
 s'_j=
 \begin{cases}
 -s_j,&0\leq j\leq K,\\
 s_j,&K\leq j\leq n.
 \end{cases}
 \]
The definitions agree at \(K\). Each increment before or at
the hit changes sign, each later increment stays the same,
so all increments remain \(\pm1\).
The new path starts at \(-x\) and ends at \(y\).
Before the first zero the original integer path is positive,
because a nearest-neighbor path cannot become negative without
visiting zero. The transformed path is therefore negative
before its first zero, and that first zero is still \(K\).

Conversely every integer nearest-neighbor path from \(-x<0\)
to \(y>0\) visits zero. Reflecting its prefix through its first
zero produces a path from \(x\) to \(y\) whose first zero is
at the same time. Applying either transformation twice is the
identity. This proves a bijection, including when the path
sets are both empty. Their count is \(N_{n,y+x}\).
Subtracting from the unrestricted count gives the positive-path
formula. At \(n=0\) there are no zero-visiting paths with
positive endpoints, and \(N_{0,y+x}=0\), so the same formula holds.
\end{proof}
\end{annexitem}

\begin{annexitem}{annex:p15.rb-ballot}{thm:p15-rb-ballot}{The strict prefix count and its weak counterpart}
\begin{proof}
Replace each \(A\) by an up-step and each \(B\) by a down-step.
All \(\binom n\alpha\) words end at \(h=\alpha-\beta\).
If \(h\leq0\), a strict lead at every nonempty prefix is
impossible at the final prefix.
Suppose \(h>0\). The first step of a strictly positive path
must go from zero to one. Reflection applied to the remaining
\(n-1\) steps counts those paths as
\[
 N_{n-1,h-1}-N_{n-1,h+1}
       =\binom{n-1}{\alpha-1}-\binom{n-1}{\alpha}.
 \]
For \(0\leq\alpha\leq n\),
\[
 \binom{n-1}{\alpha-1}=\frac\alpha n\binom n\alpha,\qquad
 \binom{n-1}{\alpha}=\frac{n-\alpha}{n}\binom n\alpha,
 \]
including the endpoints through the zero-binomial convention.
Their difference is
\((\alpha-\beta)\binom n\alpha/n\).
Division by the total word count proves the source strict formula
without ever using a factorial of a negative integer.

For a weak lead one must have \(h\geq0\).
Add one to every height: nonnegative paths from zero to \(h\)
become strictly positive paths from one to \(h+1\).
The reflection formula counts these as
\[
 N_{n,h}-N_{n,h+2}
       =\binom n\alpha-\binom n{\alpha+1}
       =\frac{\alpha-\beta+1}{\alpha+1}\binom n\alpha.
 \]
For \(h<0\) the final prefix violates the condition.
This proves the weak formula and the balanced count.
At the empty word there are no nonempty prefixes; the stated
vacuous convention explains why that case is kept outside the
ratios involving \(n\).
\end{proof}
\end{annexitem}

\begin{annexitem}{annex:p15.rb-return-counts}{thm:p15-rb-return-counts}{The telescoping positive-path count}
\begin{proof}
For \(n\geq1\), a strictly positive path at times
\(1,\ldots,2n\) ends at \(2r\) for some \(1\leq r\leq n\).
The strict ballot count, before division by the endpoint count,
gives
\[
 \mathbb P(S_1>0,\ldots,S_{2n}>0,S_{2n}=2r)
 =\frac12\{p_{2n-1,2r-1}-p_{2n-1,2r+1}\}.
 \]
Sum over \(r\). The finite sum telescopes, with all terms
beyond the reachable lattice equal to zero, to
\(\tfrac12p_{2n-1,1}\).
The binomial count shows
\[
 p_{2n-1,1}
 =2^{-(2n-1)}\binom{2n-1}{n}
 =2^{-2n}\binom{2n}{n}=u_{2n}.
 \]
Reflection of the whole path gives the same probability of
remaining strictly negative. A path with no zero at positive
times cannot change sign. Hence
\(\mathbb P(R>2n)=u_{2n}\).
For \(n=0\) this says \(1=u_0\).
Returns can occur only at even times, so for \(n\geq1\),
\[
 \mathbb P(R=2n)
 =u_{2n-2}-u_{2n}
 =u_{2n}\left\{\frac{2n}{2n-1}-1\right\}
 =\frac{u_{2n}}{2n-1}.
 \]

Condition on the first step. If it is \(1\), the remaining
return duration is the first hit of \(-1\) by the independent
future walk from zero. If it is \(-1\), the remaining duration
is the first hit of \(1\). Negating all future increments
shows that these two hitting times have the same law.
Thus \(R\) has the law of \(1+T_1\), initially as an identity
also allowing infinite times.

The proved Stirling formula gives
\[
 u_{2n}=\binom{2n}{n}4^{-n}\sim(\pi n)^{-1/2}.
 \]
In particular \(\mathbb P(R=\infty)=\lim_nu_{2n}=0\),
so \(T_1\) is finite too. The first-return mass asymptotic follows
by dividing by \(2n-1\).
For \(p>0\), the series for \(\mathbb ER^p\) has terms
\((2n)^p f_{2n}\), asymptotic to a positive constant times
\(n^{p-3/2}\). The positive series converges exactly when
\(p-3/2<-1\), or \(p<1/2\).
Replacing \(2n\) by \(2n-1\) gives the same classification
for \(T_1\). This proves the infinite means as well.
All distributional calculations in this proof came from path
counts and fixed-time independence, not from optional stopping.
\end{proof}
\end{annexitem}

\begin{annexitem}{annex:p15.rb-renewal}{thm:p15-rb-renewal}{The actual first-return decomposition}
\begin{proof}
Let \(\mathcal F_m=\sigma(X_0,\ldots,X_m)\).
For the given kernel define
\[
 h_0(x)=\mathbf1_{\{x=a\}},\qquad
 h_{r+1}(x)=\int_E h_r(y)P(x,dy).
 \]
These functions are measurable and take values in \([0,1]\).
To see measurability of each kernel integral, first check it for
indicator functions by the kernel definition, then for nonnegative
simple functions by finite linearity, and finally for bounded
nonnegative functions by increasing simple approximation and
monotone convergence.
The Markov property extends from indicators to bounded measurable
functions by the same approximation, with linearity for signed
functions. Iterating the tower identity therefore gives
\[
 \mathbb P(X_{m+r}=a\mid\mathcal F_m)=h_r(X_m).
 \]
With \(X_0=a\), taking expectations gives \(h_r(a)=v_r\).

For \(n\geq1\), if \(X_n=a\), its first positive return has
some value \(1\leq m\leq n\). The event
\(\{\tau_a^+=m\}\) belongs to \(\mathcal F_m\) and implies \(X_m=a\).
Consequently
\[
 \begin{split}
 \mathbb P_a(\tau_a^+=m,X_n=a)
 &=\mathbb E_a\{\mathbf1_{\{\tau_a^+=m\}}h_{n-m}(X_m)\}\\
 &=f_m v_{n-m}.
 \end{split}
 \]
Summing the disjoint events proves the renewal identity.
No unproved assertion about an arbitrary process restarting
at a return time has been used.

For \(0\leq s<1\), \(v_n\leq1\) gives
\(V(s)\leq(1-s)^{-1}\), and
\(\sum_{n\geq1}f_n\leq1\) gives \(F(s)\leq s\).
Multiply the renewal identity by \(s^n\) and sum.
Tonelli, or absolute convergence for complex \(|s|<1\),
gives
\[
 V(s)-1
 =\sum_{m\geq1}\sum_{j\geq0}f_m v_j s^{m+j}
 =F(s)V(s).
 \]
Because \(|F(s)|\leq F(|s|)<1\), the division is legitimate.
Let \(s\uparrow1\) in the real nonnegative series using
monotone convergence. The two alternatives at \(s=1\) follow,
including the infinite value when \(F(1)=1\).

The fair walk is a homogeneous Markov chain on the integers,
because its next independent fair step depends on the past
only through its present position.
Its odd \(u_n\)'s vanish. The binomial series at exponent
\(-1/2\) gives
\[
 \sum_{n\geq0}u_ns^n
 =\sum_{k\geq0}\binom{2k}{k}(s^2/4)^k
 =(1-s^2)^{-1/2}.
 \]
Solving the renewal identity for \(F(s)\) gives the stated
return generating function. This derives it through the
renewal count independently of the earlier exponential-martingale route.
\end{proof}
\end{annexitem}

\begin{annexitem}{annex:p15.rb-last-zero}{thm:p15-rb-last-zero}{A zero followed by a zero-free suffix}
\begin{proof}
The last zero by \(2n\) occurs at an even time.
For \(0\leq k\leq n\),
\[
 \{L_{2n}=2k\}
 =\{S_{2k}=0,\ S_{2k+1}\ne0,\ldots,S_{2n}\ne0\}.
 \]
The future increments are independent of the first \(2k\)
increments. On \(S_{2k}=0\), their zero-free condition has
probability \(u_{2n-2k}\) by the proved no-return identity.
Multiplication gives the formula. When \(k=n\), the suffix
is empty and has probability one, exactly \(u_0\).
When \(k=0\), the initial zero has probability one.
Every path has exactly one last zero, so these probabilities
sum to one. Their symmetry is immediate from the product.
\end{proof}
\end{annexitem}

\begin{annexitem}{annex:p15.rb-arcsine-limit}{thm:p15-rb-arcsine-limit}{The source step density and the singular endpoints}
\begin{proof}
Fix \(0<a<b<1\). Stirling's asymptotic for \(u_{2m}\)
implies, uniformly for \(k/n\) in any fixed compact
subinterval of \((0,1)\),
\[
 n\,u_{2k}u_{2n-2k}
   -\frac1{\pi\sqrt{(k/n)(1-k/n)}}\longrightarrow0.
 \]
Indeed both \(k\) and \(n-k\) tend to infinity uniformly
there, and the supremum of the asymptotic relative error
over integers \(m\geq c n\) tends to zero for each \(c>0\).
The remaining continuous expression is uniformly bounded
on the chosen compact interval.

Set \(a_n=\lceil na\rceil/n\), \(b_n=\lfloor nb\rfloor/n\),
and, for \(0\leq k\leq n\), define the step function
\[
 h_n(x)=n\,\mathbb P(L_{2n}=2k)
                       \quad\text{on }[k/n,(k+1)/n).
 \]
Then the exact finite sum is
\[
 \mathbb P(a\leq L_{2n}/(2n)\leq b)
 =\sum_{k=\lceil na\rceil}^{\lfloor nb\rfloor}
              \mathbb P(L_{2n}=2k)
 =\int_{a_n}^{b_n+1/n}h_n(x)\,dx.
 \]
For all sufficiently large \(n\), these integration intervals
lie in one compact subset of \((0,1)\).
The preceding lattice estimate and uniform continuity of
\(h(x)=1/(\pi\sqrt{x(1-x)})\) there give uniform convergence
of \(h_n\) to \(h\) there.
Thus the integral difference is bounded by interval length
times the uniform error, plus the integral of \(h\) over the
two vanishing endpoint adjustments. This tends to zero,
proving the source interval limit with the correct \(2k\)
mass and rounded indices.

For normalization put \(x=y^2\), then \(y=\sin\theta\)
on \([0,\pi/2]\). The one-variable substitution gives
\[
 \int_a^b h(x)\,dx
 =\frac2\pi\int_{\sqrt a}^{\sqrt b}(1-y^2)^{-1/2}\,dy
 =\frac2\pi\{\arcsin\sqrt b-\arcsin\sqrt a\}.
 \]
Taking \(a\downarrow0,b\uparrow1\) gives total integral one.
These integrals also show that the displayed CDF is continuous
at both endpoints and has no endpoint atoms.

To obtain full weak convergence rather than only interior
interval limits, fix \(0<\delta<1/2\). The proved result gives
probability in \([\delta,1-\delta]\) tending to
\(\int_\delta^{1-\delta}h\), which tends to one as
\(\delta\downarrow0\). The probability outside that compact
interval can therefore be made uniformly small in the limit.
For any \(0<x<1\), use the limit on \([\delta,x]\)
and the small mass below \(\delta\), then let
\(\delta\downarrow0\). This proves convergence of the CDF
at \(x\). The endpoint masses are at most the corresponding
small endpoint neighborhoods, giving convergence at zero and one.
Outside \([0,1]\) the CDFs already have the required values.
No uniform local approximation at the singular endpoints was assumed.
\end{proof}
\end{annexitem}

\begin{annexitem}{annex:p15.rb-occupation}{thm:p15-rb-occupation}{The separate first-excursion induction}
\begin{proof}
Between two successive zeros, all interior vertices have the
same sign; otherwise a nearest-neighbor integer path would visit
zero earlier. Every complete excursion has even length.
The final incomplete excursion before time \(2n\) also has
even length, because it begins at an even zero time.
Each positive excursion contributes all its edges to \(A_{2n}\),
and each negative excursion contributes none. Hence \(A_{2n}\)
is even, including when the final segment of the path is incomplete.

Write \(\beta_{k,n}=\mathbb P(A_{2n}=2k)\), for \(0\leq k\leq n\),
and \(f_{2m}=\mathbb P(R=2m)\), \(m\geq1\), in this walk.
For \(n=0\), \(\beta_{0,0}=1=u_0^2\).
For the upper endpoint \(k=n\), all vertices are nonnegative.
The strict-positive survival probability at horizon \(2n\),
for \(n\geq1\), is
\[
 \begin{split}
 \tfrac12u_{2n}
 &=\mathbb P(S_1>0,\ldots,S_{2n}>0)\\
 &=\tfrac12\,\mathbb P(S_1\geq0,\ldots,S_{2n-1}\geq0)\\
 &=\tfrac12\,\mathbb P(S_1\geq0,\ldots,S_{2n}\geq0)
 =\tfrac12\beta_{n,n}.
 \end{split}
 \]
The second equality conditions on the first up-step and translates
the independent remaining walk down by one. The third uses parity:
at odd time \(2n-1\), a nonnegative integer height is at least
one, so its next step remains nonnegative.
Thus \(\beta_{n,n}=u_{2n}\).
Negating every increment proves \(\beta_{0,n}=u_{2n}\).
This includes the source's \(n=1\) endpoint base cases.

Proceed by induction on \(n\). For \(1\leq k\leq n-1\),
the path must spend edges on both sides of zero, so its first
return is \(R=2m<2n\). Its first excursion is either positive
and contributes \(2m\) edges, or negative and contributes zero.
For fixed \(m\), negating the first excursion bijects its
positive and negative possibilities, each with probability
\(f_{2m}/2\). The remaining fixed-time increments are independent
of that first-return event. Therefore
\[
 \beta_{k,n}
 =\frac12\sum_{m=1}^k f_{2m}\beta_{k-m,n-m}
  +\frac12\sum_{m=1}^{n-k} f_{2m}\beta_{k,n-m}.
 \]
Every shorter-horizon term is covered by the induction
hypothesis, including its endpoint cases.
Substitution gives
\[
 \beta_{k,n}
 =\frac12u_{2n-2k}\sum_{m=1}^k f_{2m}u_{2k-2m}
  +\frac12u_{2k}\sum_{m=1}^{n-k}f_{2m}u_{2n-2k-2m}.
 \]
The two sums are \(u_{2k}\) and \(u_{2n-2k}\), respectively,
by the first-return renewal decomposition.
Their substitution gives \(\beta_{k,n}=u_{2k}u_{2n-2k}\).
This closes the induction. The law agrees with the last-zero
law, so the proved arcsine limit transfers directly.
\end{proof}
\end{annexitem}

\begin{annexitem}{annex:p15.rb-laplace-premise}{thm:p15-rb-laplace-premise}{Constructing the positive stable law and proving the transform criterion}
\begin{proof}
The density is nonnegative. Substitute \(u=y^{-1/2}\)
in its integral on \(y>0\). Then
\[
 \int_0^\infty g(y)\,dy
       =\sqrt{\frac2\pi}\int_0^\infty e^{-u^2/2}\,du=1
 \]
by the proved Gaussian normalization.
For \(a\geq0\), put
\[
 J(a)=\int_0^\infty
              \exp\{-u^2/2-a^2/(2u^2)\}\,du.
 \]
It is finite and \(J(0)=\sqrt{\pi/2}\).
On a compact interval of positive \(a\)'s, the derivative
integrand is dominated by a constant times
\[
 u^{-2}e^{-u^2/2-c/u^2}
 \]
for some \(c>0\). This is integrable: for \(u\geq1\) a
Gaussian tail dominates it, while for \(0<u\leq1\) the
substitution \(v=1/u\) gives a Gaussian-tail integral.
Thus differentiation under the integral is justified, and
the positive real substitution \(v=a/u\) gives
\[
 J'(a)=-\int_0^\infty \frac a{u^2}
                 e^{-u^2/2-a^2/(2u^2)}\,du=-J(a).
 \]
Hence \(e^aJ(a)\) is constant on \((0,\infty)\).
Dominated convergence as \(a\downarrow0\), with majorant
\(e^{-u^2/2}\), identifies the constant:
\(J(a)=\sqrt{\pi/2}e^{-a}\).
The same \(u=y^{-1/2}\) substitution now gives
\[
 \mathbb Ee^{-\lambda Y}
    =\sqrt{\frac2\pi}J(\sqrt{2\lambda})
    =e^{-\sqrt{2\lambda}},
 \]
including \(\lambda=0\).

For uniqueness of Laplace transforms, if finite nonnegative
variables have identical transforms, their transforms
\(U=e^{-Z}\in(0,1]\) have identical integer moments.
Probability laws on \([0,1]\) are determined by their moments,
by the proved compact-support moment theorem (or its Bernstein
polynomial approximation proof). Hence the \(U\)-laws agree.
The equivalence \(\{Z\leq x\}=\{U\geq e^{-x}\}\) for \(x\geq0\)
then identifies the original laws.

For convergence, let \(U_n=e^{-Z_n}\) and \(U=e^{-Y}\).
The assumed Laplace limits at nonnegative integer arguments give
convergence of every moment of \(U_n\) to that of \(U\).
The proved method of moments, with the compact determinate
target law, gives \(U_n\Rightarrow U\).
At every continuity point \(x\geq0\) of the law of \(Y\),
\(\mathbb P(U=e^{-x})=\mathbb P(Y=x)=0\).
Weak convergence of \(U_n\) at that threshold thus gives
\[
 \mathbb P(Z_n\leq x)=\mathbb P(U_n\geq e^{-x})
                   \longrightarrow\mathbb P(U\geq e^{-x})
                   =\mathbb P(Y\leq x).
 \]
At negative \(x\) all these probabilities are zero.
This proves weak convergence without asserting continuity of
the logarithm at zero or assuming a missing Laplace convergence theorem.

Finally independent copies give
\[
 \mathbb E\exp\{-\lambda(Y_1+\cdots+Y_m)/m^2\}
    =\{e^{-\sqrt{2\lambda/m^2}}\}^m
    =e^{-\sqrt{2\lambda}}.
 \]
Laplace uniqueness proves the strict stability statement,
with index \(1/2\) because the sum scale is \(m^{1/(1/2)}=m^2\).
\end{proof}
\end{annexitem}

\begin{annexitem}{annex:p15.rb-return-scaling}{thm:p15-rb-return-scaling}{The exact return shift and the square-scale limit}
\begin{proof}
The first-return tail proves \(R<\infty\) almost surely, and
\(R\stackrel{\mathrm d}=1+T_1\) proves the same for \(T_1\).
Each \(R_m\) is a stopping time: the event \(R_m\leq n\)
means at least \(m\) zeros among \(S_1,\ldots,S_n\).
Each level hit \(T_m\) is a stopping time as well.
The independent-increment restart lemma
(Theorem~\ref{thm:p15-os-independent-restart}) shows inductively
that the successive return durations are independent copies of \(R\),
and the successive upward one-level durations are independent
copies of \(T_1\). Each is finite, so all times in question
are finite almost surely.
Consequently
\[
 R_m\stackrel{\mathrm d}=R^{(1)}+\cdots+R^{(m)}
 \stackrel{\mathrm d}=m+T_1^{(1)}+\cdots+T_1^{(m)}
 \stackrel{\mathrm d}=m+T_m.
 \]
These are equalities of laws, not an assertion of equality
along a common walk.

Write \(F(s)=1-\sqrt{1-s^2}\), the combinatorially derived
return generating function. For \(\lambda>0\),
\[
 \mathbb Ee^{-\lambda R_m/m^2}
 =\left\{1-\sqrt{1-e^{-2\lambda/m^2}}\right\}^{m}.
 \]
Set \(b_m=\sqrt{1-e^{-2\lambda/m^2}}\).
The elementary exponential expansion gives
\(b_m\to0\) and \(mb_m\to\sqrt{2\lambda}\).
The logarithm bound
\(|\log(1-b_m)+b_m|\leq b_m^2\), valid eventually,
has \(mb_m^2\to0\). Thus the transform tends to
\(e^{-\sqrt{2\lambda}}\). At zero it equals one.
The proved nonnegative-transform criterion yields
\(R_m/m^2\Rightarrow Y\).

For \(0<s<1\), \(R\stackrel{\mathrm d}=1+T_1\) gives
\(\mathbb Es^{T_1}=F(s)/s\).
Therefore
\[
 \mathbb Ee^{-\lambda T_m/m^2}
 =e^{\lambda/m}
    \left\{1-\sqrt{1-e^{-2\lambda/m^2}}\right\}^{m}
 \longrightarrow e^{-\sqrt{2\lambda}}.
 \]
The same criterion proves the level-hit limit and its exact
normalization constant. No general stable domain-of-attraction
theorem or Brownian first-passage theorem has been presumed.
\end{proof}
\end{annexitem}

\begin{annexitem}{annex:p15.rb-calculations}{thm:p15-rb-calculations}{Worked path counts and hypothesis boundaries}
\begin{proof}
For length four from one to one, the unrestricted count is
\(N_{4,0}=\binom42=6\), and the zero-hitting count is
\(N_{4,2}=\binom43=4\).
The two positive paths are
\((1,2,3,2,1)\) and \((1,2,1,2,1)\).
For a concrete reflection pair, \((1,2,1,0,1)\) first hits
zero at time three and maps to \((-1,-2,-1,0,1)\).
For three \(A\)'s and two \(B\)'s, the strict words are
\(AAABB,AABAB\). The weak words are those two together with
\(AABBA,ABAAB,ABABA\).
They give the stated two and five of ten counts.

Since \(u_0=1,u_2=1/2,u_4=3/8\), the last-zero and
occupation formulas give masses \(3/8,1/4,3/8\).
The two displayed paths have the stated last zeros, positive-side
edges and strictly positive vertices by direct inspection of
their four listed increments.

To compute moments, put \(w_j=u_{2j}=\binom{2j}{j}4^{-j}\)
and \(W(z)=\sum_{j\geq0}w_jz^j=(1-z)^{-1/2}\), \(|z|<1\).
Termwise differentiation is valid on compact subintervals of
\((-1,1)\): since \(0\leq w_j\leq1\), the derivative series
are dominated by the convergent series
\(\sum j r^{j-1}\), \(\sum j(j-1)r^{j-2}\) for each \(0<r<1\).
The coefficient of \(z^n\) in \(zW'(z)W(z)\) is
\(\sum_{k=0}^n k w_k w_{n-k}\), while
\[
 zW'(z)W(z)=\tfrac12z(1-z)^{-2}.
 \]
The coefficient is \(n/2\).
Similarly
\[
 z^2W''(z)W(z)=\tfrac34z^2(1-z)^{-3}
 \]
has coefficient \(3n(n-1)/8\), giving
\(\mathbb E[K(K-1)]=3n(n-1)/8\).
Add \(\mathbb EK=n/2\) and subtract \(n^2/4\) to obtain
\(\operatorname{Var}(K)=n(n+1)/8\).
At \(n=1\), the same formulas hold with the second factorial
moment zero. The assertions about the normalized fractions follow.

For the finite mass profile use
\(w_{j+1}/w_j=(2j+1)/(2j+2)\).
For \(0\leq k<n\), the ratio of adjacent masses is
\[
 \frac{w_{k+1}w_{n-k-1}}{w_kw_{n-k}}
 =\frac{2k+1}{2k+2}\,
                  \frac{2(n-k)}{2(n-k)-1}.
 \]
It is at most one exactly when \(k\leq(n-1)/2\).
Together with symmetry this proves endpoint maxima and
one or two central minima, including the admissible-lattice
qualification when a perfectly equal division is impossible.

The arcsine CDF has \(H(1/4)=1/3\), \(H(3/4)=2/3\),
and \(H(1/2)=1/2\).
This gives the middle/outer probabilities.
No zero after time \(n\) is exactly \(L_{2n}\leq n\);
the weak limit at \(1/2\) proves its probability tends to \(1/2\).
It does not give an exact finite-horizon calendar statement.

For biased increments, the path \((1,0,1)\) has step word
\((-1,+1)\) and probability \(qp\).
Its reflected path \((-1,0,1)\) has \((+1,+1)\) and
probability \(p^2\). The combinatorial bijection remains
valid, but the equality of these probabilities requires \(p=q\)
when \(0<p<1\).

For a process without the stated homogeneous Markov premise,
take the deterministic sequence
\[
 X_0=a,\ X_1=b,\ X_2=a,\ X_3=b,\ X_n=b\quad(n\geq4),
 \qquad a\ne b.
 \]
It has \(f_2=1\), \(v_2=1\), and \(v_4=0\).
The proposed renewal sum at \(n=4\) would equal
\(f_2v_2=1\), a contradiction. This does not contradict the
proved theorem, because the transitions from \(b\) differ
between its first and later visits.
Finally in the homogeneous chain that stays at \(a\) forever,
\(\tau_a^+=1\), \(f_1=1\), and \(v_n=1\), giving
\(V(s)=1/(1-s)\) correctly.
If the initial hit at time zero were substituted for the first
positive return, all \(f_n\), \(n\geq1\), would be zero and
the renewal identity would fail already at \(n=1\).
\end{proof}
\end{annexitem}
